\section{The real counting branch, including its edge roots}\label{sec:counting}

The exact moment formula is now available. We locate its physical term by a real root system. The amplitude calculation will require an inverse for the counting Jacobian that is uniform in the number of roots and in the twist, including at the two extreme roots.

At \(l=\pi/8\), take the real slice \(t_j=iy_j\), \(1\le j\le N\), and homogeneous sites zero. Define
\[
 p_k(y)=\frac{\sin(2kl)}{\pi(\cosh2y-\cos2kl)},\qquad
 p=p_1,\qquad r=p_1-p_2,\qquad
 R(y)=\int_{-\infty}^y r(v)\,dv,
\]
and, for an ordered configuration \(y_1<\cdots<y_N\), define
\[
 T(y)=N\int_{-\infty}^y p(v)\,dv+\sum_{j=1}^N R(y-y_j).
\]
The phases in \(D_i\) from the auxiliary residue calculation satisfy
\[
 (2\pi i)^{-1}\partial_y\log([iy+k]/[iy-k])=p_k(y)
\]
for integer shifts \(k\), which gives these phase derivatives. We include every pair term in \(T(y)\), including the term at the same atom when \(y=y_i\). The self phase ratio is one; its two derivatives will cancel in the Jacobian below.

With the convention \(\widehat f(t)=\int_{\mathbb R}e^{ity}f(y)\,dy\), substitution \(v=e^{2y}\), partial fractions, and the Beta integral give
\[
 \widehat p_k(t)=\frac{\sinh((4-k)lt)}{\sinh(4lt)}
 \quad(0<k<4).
\]
In particular, \(\int p=3/4\) and \(\int r=1/4\). Set
\[
 \widehat\rho(t)=\frac1{2\cosh(lt)-1}
 =\frac{\sinh(2lt)+\sinh(lt)}{\sinh(3lt)}.
\]
Inverting the transform gives the explicit positive density
\[
 \rho(y)=\frac{2\sqrt3}{3\pi}
 \left(\frac1{\cosh(8y/3)-1/2}
       +\frac1{\cosh(8y/3)+1/2}\right).
\]
It has mass one and solves \(\rho=p+r*\rho\). Write
\[
 \alpha_{\rm r}=\frac83,\qquad
 x(y)=N\int_{-\infty}^y\rho(v)\,dv,\qquad
 y_N(u)=x^{-1}(u),\qquad a_N(u)=N\rho(y_N(u)).
\]
The displayed density gives \(\rho(y)\sim c_\rho e^{-\alpha_{\rm r}|y|}\) at both ends and \(|\rho^{(m)}|\le C_m\rho\) for every fixed \(m\). Consequently
\begin{equation}\label{eq:quantile-density}
 a_N(u)\asymp\min(u,N-u)\quad(0<u<N),\qquad
 \frac{a_N(u)}{\alpha_{\rm r}u}\longrightarrow1
 \quad\hbox{as }u/N\longrightarrow0,
\end{equation}
with the analogous assertion at \(N-u\). These statements are uniform in \(N\). For fixed positive \(u,v\),
\[
 y_N(u)-y_N(v)\longrightarrow
 \alpha_{\rm r}^{-1}\log(u/v)\qquad(N\longrightarrow\infty).
\]

Let
\[
 E(y)=\#\{j:y_j\le y\}-x(y).
\]
Integration by parts, using \(\rho=p+r*\rho\), gives \(T=x+r*E\). We solve
\[
 T(y_i)=i-\tfrac12+\delta,\qquad |\delta|\le\tfrac14,
 \qquad s=\frac{\delta}{1-\int r}=\frac{4\delta}{3}.
\]
In quantile coordinates \(x_i=x(y_i)\), let
\[
 \mathcal B_{N,s}=
 \left\{\boldsymbol x:
 |x_i-(i-\tfrac12+s)|\le0.21,\quad
 0.09\le x_i\le N-0.09,\quad 1\le i\le N\right\}.
\]
This is a nonempty product rectangle for \(N\ge2\); its coordinate intervals are separated by at least \(0.58\), so its points give ordered rapidities. For a vector \(\boldsymbol x\) in this rectangle, \(T_{\boldsymbol x}\) denotes the function with atoms \(y_j=y_N(x_j)\), and set
\[
 F_i(\boldsymbol x,\delta)=T_{\boldsymbol x}(y_N(x_i))
          -(i-\tfrac12+\delta),\qquad
 a_i=N\rho(y_i),\qquad
 f_i=\left.\partial_yT_{\boldsymbol x}(y)\right|_{y=y_i}
       =a_i+(r'*E)(y_i).
\]
The derivative defining \(f_i\) holds the atoms fixed. The chain rule gives the actual Jacobian in quantile coordinates:
\begin{equation}\label{eq:quantile-jacobian}
 H_{ij}:=\frac{\partial F_i}{\partial x_j}
 =\delta_{ij}\frac{f_i}{a_i}
       -\frac{r(y_i-y_j)}{a_j}.
\end{equation}
In particular the self contribution \(r(0)/a_i\) in \(f_i/a_i\) is subtracted on the diagonal.

\begin{proposition}[The real counting branch]\label{prop:real-counting}
There are absolute constants \(N_0,C\) such that, for every \(N\ge N_0\) and \(|\delta|\le1/4\), the consecutive equations
\[
 T(y_i)=i-\tfrac12+\delta,\qquad 1\le i\le N,
\]
have exactly one ordered real solution. Its quantiles belong to \(\mathcal B_{N,s}\). Throughout \(\mathcal B_{N,s}\), the Jacobian \(H\) in \eqref{eq:quantile-jacobian} is invertible and \(\|H^{-1}\|_{\ell^\infty\to\ell^\infty}\le C\). For each fixed \(N\), the solution depends locally analytically on \(\delta\) and admits local analytic continuation under small analytic perturbations of the equations.
\end{proposition}

\begin{proof}
We first give the boundary signs and the Jacobian estimates. The elementary kernel calculation at the end of the proof will supply
\begin{equation}\label{eq:counting-kernel-numbers}
 \|r\|_1<0.295,\qquad
 R(y)>-0.012,\qquad
 \int_{\alpha_{\rm r}^{-1}\log(200/27)}^\infty |r(y)|\,dy<0.012.
\end{equation}

\paragraph{Boundary signs and containment.}
The quantile intervals imply
\[
 \|E+s\|_\infty\le0.71.
\]
Indeed, between two successive atoms \(E+s=j-x+s\), and the two neighboring interval bounds put this value between \(-0.71\) and \(0.71\); the two exterior intervals obey the same bound. Since \(\delta=(1-\int r)s\), the residual has the useful form
\[
 F_i=x_i-(i-\tfrac12+s)+(r*(E+s))(y_i).
\]
On an untrimmed upper face \(x_i=i-\tfrac12+s+0.21\) it is positive, and on the corresponding lower face it is negative, because
\[
 \bigl|(r*(E+s))(y_i)\bigr|
 \le0.71\|r\|_1<0.20945<0.21.
\]
This also gives containment before the two trims. For any ordered solution of the consecutive equations, put
\(M=\max_i|x_i-(i-\tfrac12+s)|\). The same interval argument now gives
\(\|E+s\|_\infty\le\tfrac12+M\), including the exterior intervals since
\(|s|\le1/3\). Therefore
\[
 M\le0.295(\tfrac12+M),\qquad
 M\le\frac{0.1475}{0.705}<0.21.
\]

Only the lower face for the first quantile and the upper face for the last quantile can be trimmed. We prove a slightly stronger statement for the first one: throughout the untrimmed box, \(x_1\le0.09\) forces \(T(y_1)<1/4\) for all sufficiently large \(N\), uniformly in \(s\). Put \(b=1+s\in[2/3,4/3]\). All atoms \(x_j\), \(j\ge2\), lie above \(b\), and
\[
 \eta_N(u)=\#\{j\ge2:x_j\le u\}-(u-b),\qquad b\le u\le N,
\]
satisfies \(\eta_N(b)=0\) and \(\|\eta_N\|_\infty\le0.71\). Comparing the atoms with the continuum identity for \(x(y_1)\), and separating the first atom, gives
\[
 T(y_1)=x_1+R(0)-\int_0^b R(y_1-y_N(u))\,du
       +\int_b^N R(y_1-y_N(u))\,d\eta_N(u).
\]
Here \(R(0)=1/8\), since \(r\) is even and has mass \(1/4\). The boundary terms in integration by parts in the last integral vanish: \(\eta_N(b)=0\), and \(R(y_1-y_N(u))\to0\) as \(u\uparrow N\). It follows that
\[
 T(y_1)\le x_1+\tfrac18+0.012b+
 0.71\int_{y_N(b)-y_1}^\infty |r(v)|\,dv.
\]
The exponential tail of \(\rho\) gives
\[
 e^{\alpha_{\rm r}(y_N(b)-y_N(x_1))}
   =\frac b{x_1}(1+o(1))
\]
uniformly for \(0<x_1\le0.09\) and \(2/3\le b\le4/3\): both inverse quantiles are in the left tail, where the relative tail asymptotic is uniform. Thus the last integral is at most \(0.012+o(1)\) by \eqref{eq:counting-kernel-numbers}; this includes \(x_1\to0\), when the lower limit tends to infinity. We obtain the strict bound
\[
 T(y_1)\le0.09+0.125+\tfrac43(0.012)+0.71(0.012)+o(1)
 =0.23952+o(1)<\tfrac14.
\]
This proves the outward sign at \(x_1=0.09\) and rules out \(x_1\le0.09\) for a consecutive solution. Reflection sends \(y_i\) to \(-y_{N+1-i}\), \(s\) to \(-s\), and \(T(y)\) to \(N-T(-y)\). It gives the strict outward sign at \(x_N=N-0.09\) and rules out \(x_N\ge N-0.09\) for a solution.

\paragraph{The Jacobian in the bulk.}
Write \(d_i=f_i/a_i\) and
\[
 S_i=\sum_{j=1}^N\frac{|r(y_i-y_j)|}{a_j}.
\]
We will prove that there are absolute \(\epsilon>0\) and \(N_0\) such that
\begin{equation}\label{eq:counting-dominance}
 d_i-S_i\ge\epsilon
 \quad\hbox{for every }\boldsymbol x\in\mathcal B_{N,s},
 \quad N\ge N_0,\quad |\delta|\le\tfrac14.
\end{equation}
Since \(r(0)>0\), this is exactly a lower bound for
\(H_{ii}-\sum_{j\ne i}|H_{ij}|\). Also \(d_i\) is uniformly bounded above: \(r'*E=r'*(E+s)\) is bounded by \(0.71\|r'\|_1\), while \eqref{eq:quantile-density} and the trims bound \(a_i\) away from zero.

We record the uniform control of the kernel sums used here and below. The quantile intervals
\([x_j-0.04,x_j+0.04]\) are disjoint and lie in \((0,N)\). On each such interval \(a_N(u)\asymp a_j\), by \eqref{eq:quantile-density}; their rapidity lengths are therefore comparable to \(1/a_j\), and their rapidities stay a bounded distance from \(y_j\). The exponential decay of \(r\) and its derivatives then implies, for each of these kernels \(K\),
\begin{equation}\label{eq:counting-kernel-tails}
 \sup_{\substack{N\ge2,\ |s|\le1/3,\ \boldsymbol x\in\mathcal B_{N,s}\\1\le i\le N}}
 \sum_{|y_j-y_i|>L}\frac{|K(y_i-y_j)|}{a_j}
 \longrightarrow0\qquad(L\longrightarrow\infty).
\end{equation}
For example, compare each summand with the integral of a constant multiple of \(e^{-c|y_i-z|}\) over its rapidity interval; the intervals are disjoint. This also bounds the full sums uniformly.

Consider a sequence of rows for which \(\min(x_i,N-x_i)\to\infty\). Then \(a_i\to\infty\), so \(d_i=1+o(1)\). On any fixed rapidity window about \(y_i\), boundedness of \(\rho'/\rho\) gives \(N\rho\asymp a_i\to\infty\) uniformly. In quantile coordinates \(a_N'(u)=\rho'(y_N(u))/\rho(y_N(u))\) is bounded. Thus on the unit quantile cells containing the atoms, \(a_N\) changes by \(O(1)\), its relative change tends to zero, and the rapidity mesh tends to zero. Sampling at \(x_j\), whose offset in each cell is bounded by the box, gives
\[
 \sum_j\frac{|r(y_i-y_j)|}{a_j}\longrightarrow\int_{\mathbb R}|r(y)|\,dy.
\]
The passage from a fixed window to the full line follows from
\eqref{eq:counting-kernel-tails}. Hence the limiting bulk margin is
\(1-\|r\|_1>0.705\).

\paragraph{The two limiting edge deficits.}
It remains to examine sequences for which \(\min(x_i,N-x_i)\) stays bounded. Reflecting if necessary, take a subsequence with \(x_i\to x\ge0.09\), \(s\to s_*\), and every fixed left quantile \(x_j\to v_j\). The index \(i\) is fixed on this subsequence. Set
\[
 b=1+s_*\in[2/3,4/3],\qquad Q=\frac bx,\qquad
 h(u)=r\bigl(\alpha_{\rm r}^{-1}\log u\bigr)\quad(u>0).
\]
Write \(h_+=\max(h,0)\) and \(h_-=\max(-h,0)\).
Let \(\mu=\sum_{j\ge1}\delta_{v_j}\). The discrepancy
\[
 D(v)=\mu((0,v])-v+s_*
\]
has absolute value at most \(0.71\). The tail quantile limits and
\eqref{eq:counting-kernel-tails} give
\begin{equation}\label{eq:limiting-edge-margin}
 \lim\alpha_{\rm r}(d_i-S_i)
 =\alpha_{\rm r}
 +\frac1x\int_0^\infty h(v/x)(\mu(dv)-dv)
 -\sum_{j\ge1}\frac{|h(v_j/x)|}{v_j}.
\end{equation}
The signed integral here is a paired difference, defined without subtracting two divergent integrals by
\[
 \int_0^\infty h(v/x)(\mu(dv)-dv)
 :=-\int_0^\infty D(v)\frac1x h'(v/x)\,dv.
\]
It converges because \(D\) is bounded and \(h\) has finite total variation. To justify \eqref{eq:limiting-edge-margin}, first restrict to bounded rapidity offsets, where
\(a_N(x_i)\to\alpha_{\rm r}x\) and
\(y_N(v)-y_N(x_i)\to\alpha_{\rm r}^{-1}\log(v/x)\).
Then remove the restriction using \eqref{eq:counting-kernel-tails} for the absolute sum and the bound on \(E+s\) against the exponential tails of \(r'\) for the signed correction.

Write \(\operatorname{TV}_I g\) for the total variation of \(g\) on \(I\). For the atoms \(j\ge2\), the discrepancy
\[
 D_b(v)=\#\{j\ge2:v_j\le v\}-(v-b),\qquad v\ge b,
\]
also has absolute value at most \(0.71\) and vanishes at \(b\). Integration by parts against \(D\) or \(D_b\) bounds a signed sum by \(0.71\) times the variation of its kernel.

For a nonfirst atom, the exact range inherited from its quantile interval is
\[
 i\ge2,\qquad
 b+i-1.71\le x\le b+i-1.29,\qquad
 x\ge b+0.29\ge\frac{287}{300},\qquad 0<Q<1.
\]
The signed correction in \eqref{eq:limiting-edge-margin} can reduce the baseline by at most \(0.71\operatorname{TV}(h)/x\). Separate the first atom from the absolute sum. Since \(v_1\ge0.09\), \(h_-\le0.012\), and \(h_+(u)/u<0.55\), its contribution is at most
\(\max(0.012/0.09,0.55/x)\). Comparing the other atoms with Lebesgue measure on \([b,\infty)\) gives the following upper bound for the total deficit from \(\alpha_{\rm r}\):
\begin{equation}\label{eq:nonfirst-edge-deficit}
 \begin{aligned}
 \mathcal D_{\ne1}(x,b)={}&
 \frac{0.71\operatorname{TV}(h)}x+
 \max\left(\frac{0.012}{0.09},\frac{0.55}x\right)
 +\int_Q^\infty |h(u)|\,\frac{du}{u}\\
 &\quad
 +\frac{0.71}{b}Q\operatorname{TV}_{[Q,\infty)}(|h|/u).
 \end{aligned}
\end{equation}
The last factor follows from
\(\operatorname{TV}_{[b,\infty)}(|h(v/x)|/v)
=x^{-1}\operatorname{TV}_{[Q,\infty)}(|h|/u)\).

For the first atom \(x=v_1\), the exact domain is
\[
 \max(0.09,b-0.71)\le x\le b-0.29,\qquad
 1<Q=\frac bx\le\frac{400}{27}<14.9.
\]
Its self terms in \eqref{eq:limiting-edge-margin} cancel, because
\(h(1)=r(0)>0\). Pairing the remaining signed atoms with Lebesgue measure only above \(b\) leaves the ordinary integral from \(0\) to \(b\). The deficit is therefore at most
\begin{equation}\label{eq:first-edge-deficit}
 \begin{aligned}
 \mathcal D_1(Q,b)={}&
 \int_0^Q h(u)\,du
 +\frac{0.71}{b}Q\operatorname{TV}_{[Q,\infty)}h
 +\int_Q^\infty |h(u)|\,\frac{du}{u}\\
 &\quad
 +\frac{0.71}{b}Q\operatorname{TV}_{[Q,\infty)}(|h|/u).
 \end{aligned}
\end{equation}
These are the two estimates that must be kept strictly below \(8/3\).

\paragraph{Elementary bounds for the two deficits.}
Put \(a=1/\sqrt2\), \(c=1-a\), \(z=u^{3/4}\), and
\(v=(z+z^{-1})/2=\cosh(\tfrac34\log u)\). Direct substitution gives
\begin{equation}\label{eq:edge-kernel-formulas}
 \begin{aligned}
 h(u)&=\frac1\pi\left(\frac a{v-a}-\frac1v\right)
      =\frac{a-cv}{\pi v(v-a)},\\
 \mathcal P(u):=\int_0^u h(w)\,\frac{dw}{w}
 &=\frac8{3\pi}\left(\frac\pi4+
      \arctan(z/a-1)-\arctan z\right).
 \end{aligned}
\end{equation}
In particular \(h(u)=h(1/u)\). Its positive interval is
\((\zeta^{-1},\zeta)\), where
\[
 v_0=\frac a{1-a},\qquad
 \zeta=(v_0+\sqrt{v_0^2-1})^{4/3},\qquad 7.5<\zeta<8.
\]
The maximum is \(h(1)=\sqrt2/\pi<0.451\). The two negative minima occur at
\(v=a/(1-\sqrt a)\), with common absolute value
\[
 m=\frac{(1-\sqrt a)^2}{\pi a}<0.012.
\]
Thus \(-0.012<h<0.451\) and
\[
 \operatorname{TV}(h)=2h(1)+4m<0.95.
\]
The primitive has \(\mathcal P(\infty)=2/3\) and
\(\mathcal P(u)+\mathcal P(1/u)=2/3\). Evaluation of the formula in
\eqref{eq:edge-kernel-formulas} gives
\[
 -0.0294<\mathcal P(\zeta^{-1})<-0.0290,\qquad
 0.6957<\mathcal P(200/27)<0.6961.
\]
Since the primitive decreases, then increases, then decreases on the three sign intervals of \(h\), these bounds imply
\[
 \int_0^\infty |h(u)|\,\frac{du}{u}
 =\frac23-4\mathcal P(\zeta^{-1})<0.786,\qquad
 \int_1^\infty |h(u)|\,\frac{du}{u}<0.393.
\]
They also prove \eqref{eq:counting-kernel-numbers}. Indeed \(dy=du/(\alpha_{\rm r}u)\), so
\(\|r\|_1<0.786/\alpha_{\rm r}=0.29475<0.295\) and
\(\inf R=\mathcal P(\zeta^{-1})/\alpha_{\rm r}>-0.012\). Since
\(200/27<\zeta\), the absolute logarithmic tail there equals
\[
 \frac23-2\mathcal P(\zeta^{-1})-\mathcal P(200/27)
 <\frac23+2(0.0294)-0.6957<0.032
 =0.012\alpha_{\rm r}.
\]

We next bound the variations of \(|h|/u\). Let \(G(u)=h_+(u)/u\). On the positive interval define
\(\phi(t)=\log G(e^t)\). Differentiating the rational expression above yields
\[
 \phi''(t)=\frac9{16}
 \left(-\frac{c(av-c)}{(a-cv)^2}-\frac1{v^2}
                 +\frac{av-1}{(v-a)^2}\right)<0,
 \qquad v=\cosh(3t/4).
\]
For completeness, if \(av\le1\), every displayed term is nonpositive and the first two are negative. If \(av>1\), then on the positive interval
\(0<a-cv\le c\), \(v-a\ge a\), and
\(c(v-a)^2\ge c/2>c^2\ge(a-cv)^2\); also \(av-c>av-1\).
The first negative term therefore dominates the last positive term.
The logarithm of \(G\) is strictly concave.

The following values locate its unique mode:
\[
 \begin{array}{c|ccc}
 t&-0.40&-0.38&-0.36\\ \hline
 \phi'(t)&(0.055,0.065)&(0.011,0.019)&(-0.036,-0.027).
 \end{array}
 \qquad G(e^{-0.38})<0.54.
\]
It follows that the mode \(u_*\) lies between \(e^{-0.40}\) and
\(e^{-0.36}<0.70\). The tangent bound for the concave function \(\phi\) gives
\[
 \max_{u>0}G(u)<0.54e^{0.019(0.02)}<0.55.
\]
For \(0<Q<1\), the positive contribution to the scaled variation is
\[
 Q\operatorname{TV}_{[Q,\infty)}G
 =\begin{cases}
 Q(2G(u_*)-G(Q)),&Q<u_*,\\
 QG(Q),&Q\ge u_*.
 \end{cases}
\]
The second line is at most \(h_+(Q)<0.451\). The first is at most
\(1.10Q-h_+(Q)\). It is below \(0.55\) if \(Q\le0.5\); if
\(0.5<Q<u_*\), monotonicity of \(h\) on \([0.5,1]\) and \(h(0.5)>0.24\) give the sharper bound \(1.10(0.70)-0.24=0.53\). Thus this positive contribution is always below \(0.55\).

For the negative part on the left interval \((0,\zeta^{-1})\),
\(h_-(u)/u\) decreases: \(v h_-(u)=(cv-a)/(\pi(v-a))\) increases with \(v\), while \(v\) decreases and \(uv\) increases with \(u\). Its scaled variation above \(Q\) is therefore at most \(h_-(Q)<0.012\). There is no left contribution for \(Q\ge1\). On the right interval \((\zeta,\infty)\), the variation of \(h_-\) is at most \(2m<0.024\). The product variation bound applied to \(h_-(u)/u\) gives, with \(q=\max(Q,\zeta)\),
\[
 Q\operatorname{TV}_{[q,\infty)}(h_-/u)
 \le Q\left(\frac{0.024}{q}+\frac{0.012}{q}\right)
 <0.036\min(1,Q/7.5).
\]
For \(0<Q<1\) the full scaled variation is consequently below
\(0.55+0.012+0.036/7.5=0.5668\). Since \(b\ge2/3\),
\begin{equation}\label{eq:nonfirst-variation-number}
 \frac{0.71}{b}Q\operatorname{TV}_{[Q,\infty)}(|h|/u)
 <1.065(0.5668)=0.603642<0.605.
\end{equation}

We can now close \eqref{eq:nonfirst-edge-deficit}. If \(x\ge1.1\), then
\[
 \mathcal D_{\ne1}(x,b)
 <\frac{0.675}{1.1}+0.50+0.786+0.605
 <2.505<\frac83.
\]
If \(287/300\le x<1.1\), then
\(Q\ge(2/3)/1.1=20/33>0.60\). Inversion symmetry and \(h<0.451\) give
\[
 \int_Q^\infty |h(u)|\,\frac{du}{u}
 <0.393+0.451\log(5/3).
\]
Hence in this remaining domain
\[
 \mathcal D_{\ne1}(x,b)
 <\frac{0.675}{287/300}+\frac{0.55}{287/300}
   +0.393+0.451\log(5/3)+0.605
 <2.510<\frac83.
\]

For \eqref{eq:first-edge-deficit}, recall \(1<Q\le400/27\).
The sign pattern of \(h\) gives
\[
 Q\operatorname{TV}_{[Q,\infty)}h
 \le Qh_+(Q)+2Qm.
\]
By inversion \(Qh_+(Q)=G(1/Q)<0.55\). For \(Q\ge4\), the argument \(1/Q\) is below the mode of \(G\), so this quantity decreases with \(Q\) and is at most \(4h(4)<0.22\). It follows that
\begin{equation}\label{eq:first-signed-variation}
 Q\operatorname{TV}_{[Q,\infty)}h
 <\begin{cases}
 0.55+0.096=0.646,&1<Q\le4,\\
 0.22+0.358=0.578,&4<Q\le400/27.
 \end{cases}
\end{equation}
Here the negative terms use \(2Qm<0.096\) in the first range and
\(2Qm<2(14.9)(0.012)<0.358\) in the second. Also \(G\) decreases on
\([1,\infty)\), and only the right negative interval contributes there. The previous product bound therefore gives
\begin{equation}\label{eq:first-absolute-variation}
 Q\operatorname{TV}_{[Q,\infty)}(|h|/u)
 \le h_+(Q)+0.036
 <\begin{cases}
 0.487,&1<Q<2,\\
 0.281,&2\le Q\le400/27,
 \end{cases}
\end{equation}
where the second line uses \(h_+(Q)\le h(2)<0.245\).

Here is an explicit bound for the remaining ordinary integral. The function \(h_+\) increases on \((0,1)\), decreases on \((1,\zeta)\), and vanishes outside \((\zeta^{-1},\zeta)\). The following upper endpoint values follow from \eqref{eq:edge-kernel-formulas}:
\[
 \begin{array}{c|rrrrrrrrr}
 u&0.25&0.5&0.75&1&1.5&2&2.5&3&3.5\\ \hline
 10^3h_+(u)<&55&243&401&451&359&243&163&112&78
 \end{array}
\]
\[
 \begin{array}{c|rrrrrrrr}
 u&4&4.5&5&5.5&6&6.5&7&7.5\\ \hline
 10^3h_+(u)<&55&39&28&19&13&8&4&1.
 \end{array}
\]
Use right endpoints on the four intervals of length \(0.25\) up to \(1\), and left endpoints on intervals of length \(0.5\) from \(1\) to \(8\). The corresponding upper sums are below \(0.694\) up to \(2\) and below \(1.074\) over the whole positive interval. Thus
\begin{equation}\label{eq:first-ordinary-integral}
 \int_0^Q h(u)\,du
 <\begin{cases}0.70,&1<Q<2,\\1.08,&2\le Q\le400/27.\end{cases}
\end{equation}
All the decimal bounds used in this calculation can be checked directly from the displayed formulas. For the endpoint table, bracket \(a\) between \(0.70710678\) and \(0.70710679\) by squaring, and bracket \(z=u^{3/4}\) between consecutive multiples of \(10^{-8}\) by comparing fourth powers with \(u^3\); substitution gives the stated integer ceilings, as well as \(0.24<h(0.5)<0.245\) and \(4h(4)<0.22\). The primitive bounds use \(z=v_0-\sqrt{v_0^2-1}\) at \(\zeta^{-1}\), and \(z=(200/27)^{3/4}\) at \(200/27\). After reciprocal reduction every arctangent argument has absolute value below \(0.71\), so its first 30 alternating terms bound it with the next-term remainder. The identity \(\pi/4=\arctan(1/2)+\arctan(1/3)\) gives the needed bound for \(\pi\) by the same rule. For the three values of \(\phi'\), use
\[
 \phi'(t)=-1+\tfrac34\sinh(3t/4)
 \left(-\frac c{a-cv}-\frac1v-\frac1{v-a}\right);
\]
the first ten Taylor terms for the exponentials at \(|t|\le0.40\) have error below \(10^{-9}\). These rational interval evaluations give the displayed ranges, including \(7.5<\zeta<8\), without a numerical integration.

Finally \(\int_Q^\infty |h|\,du/u<0.393\) for \(Q>1\), and
\(0.71/b\le1.065\). Combining the correlated ranges in
\eqref{eq:first-signed-variation}--\eqref{eq:first-ordinary-integral} gives a strict total in every case:
\[
 \begin{array}{ll}
 1<Q<2:
 &\mathcal D_1
 <0.70+1.065(0.646)+0.393+1.065(0.487)
 =2.299645<8/3,\\[2mm]
 2\le Q\le4:
 &\mathcal D_1
 <1.08+1.065(0.646)+0.393+1.065(0.281)
 =2.460255<8/3,\\[2mm]
 4<Q\le400/27:
 &\mathcal D_1
 <1.08+1.065(0.578)+0.393+1.065(0.281)
 =2.387835<8/3.
 \end{array}
\]
In particular, the larger variation bound for \(Q<2\) is paired with the smaller ordinary-integral bound in that same range.

\paragraph{Existence, uniqueness, and continuation.}
The bulk margin has a fixed positive limit, and every limiting edge deficit has a fixed gap below \(\alpha_{\rm r}\). The subsequence argument above therefore proves \eqref{eq:counting-dominance} uniformly: otherwise a sequence of rows with margins tending to zero would have either a bulk subsequence or one of the two edge subsequences, contradicting the corresponding strict bound. The diagonal entries of \(H\) are positive and uniformly bounded above. Choose a sufficiently small absolute \(\eta>0\). Then every derivative of
\(\boldsymbol x\mapsto\boldsymbol x-\eta F(\boldsymbol x,\delta)\) has
\(\ell^\infty\) row norm at most \(1-\eta\epsilon\). Coordinatewise projection onto \(\mathcal B_{N,s}\) is nonexpansive, so the projected map is a contraction of the rectangle. Its fixed point exists by iteration. The strict boundary signs put that point in the interior, where a projected fixed point satisfies \(F=0\).

For any two points of the rectangle, the Jacobian averaged along their segment retains positive diagonal entries and the same strict row margin. It is invertible by the maximum-coordinate estimate, which also gives
\(\|H^{-1}\|_\infty\le\epsilon^{-1}\). Applying the averaged Jacobian to the difference of two zeros proves uniqueness. The containment argument at the start places every ordered solution in this rectangle. Finally the analytic implicit function theorem, applied for fixed \(N\) to \eqref{eq:quantile-jacobian}, proves the stated local dependence and continuation.
\end{proof}
